\section{Introduction}
\label{sec:intro}

Weight quantization is now the default way to make large language models
deployable, and the field has settled into two regimes.\blfootnote{Generative
AI tools were used in the preparation of this manuscript. All technical
claims, numbers, and citations were verified by the author against the
recorded reports and public sources, and the author takes full responsibility
for the content.} Post-training
quantization (PTQ) rounds a trained model using calibration statistics, with
error-compensating variants such as GPTQ~\citep{gptq} and rotation-based
preprocessing such as QuIP~\citep{quip}, QuIP\#~\citep{quipsharp},
QuaRot~\citep{quarot} and SpinQuant~\citep{spinquant}. Quantization-aware
training (QAT) instead learns the weights with the quantizer in the loop,
usually with a straight-through estimator~\citep{ste}, sometimes with
distillation~\citep{kd}. Ternary weights---$\{-s, 0, +s\}$, roughly two bits
per weight---sit at the aggressive end of this spectrum, from early ternary
weight networks~\citep{twn} to ternary-native training at scale
\citep{bitnet158}.

A released ternary checkpoint, however, is an opaque object. It reports a bit
width and a benchmark table, but the file itself does not say how the weights
were produced, how much of the quality lives in the released codes rather than
in the quantizer, or whether the recipe is reproducible from the public base
model. This matters: the claims that make a low-bit release interesting are
exactly the claims a user cannot check.

We report an independent forensic study of one such release: Ternary
Bonsai 2 27B, a $\sim$2\,bpw ternary model built from Qwen3.8-27B that ships as
a GGUF and runs on a public inference fork~\citep{prismbonsai2,prismfork,
qwen38}. We do not run the training recipe. We treat the released artifact as
the specimen and the public base model as the control, and we ask four
questions:

\begin{itemize}
  \item \textbf{Q1 (format).} What exactly is the storage format, and can the
  released codes be reproduced from the public base model without access to the
  training code?
  \item \textbf{Q2 (quality).} How much of the model's quality lives in the
  released weights, as opposed to the container or the quantizer?
  \item \textbf{Q3 (mechanism).} Is the residual between a public-base round-to-
  nearest solution and the released codes an error-compensation trick, or
  trained weight movement?
  \item \textbf{Q4 (evaluation).} Does a widely cited public PTQ baseline
  reproduce its published quality when calibration and evaluation are
  separated?
\end{itemize}

\paragraph{Contributions.}
\begin{enumerate}
  \item \textbf{Format and basis recovery} (\S\ref{sec:format}). From the
  released file alone, we recover the blockwise Walsh--Hadamard rotation basis,
  the axis and sign conventions, and the tensor layouts, and reproduce the
  released trits from the public base with plain absmean round-to-nearest at
  $0.920$ mean / $0.885$ minimum per-tensor agreement (chance $0.33$). The
  recovery is later confirmed exactly against metadata the model authors
  published independently: block size, transform, axis, explicit signs, and
  the inverse-rotated embedding all match.
  \item \textbf{The residual is the model} (\S\ref{sec:residual}). A
  byte-controlled, in-place replacement of all 402 ternary payloads with
  rotate-and-round weights of the public base---metadata, layouts, and
  exemptions kept byte-identical---collapses held-out perplexity from
  $18.5851$ to $23{,}606$ ($1{,}270\times$) in the same binary, corpus, and
  settings, even though the replaced codes still agree with the released ones
  on 90--95\% of weights.
  \item \textbf{The residual is trained weight movement} (\S\ref{sec:mechanism}).
  With real activation Hessians captured on a four-layer prefix of the 27B, we
  show that every error-compensated assignment we tried---GPTQ damping sweeps,
  act-order, rotation-aware Hessians, group scale search, stochastic
  calibration---moves codes \emph{away} from the released ones, and that the
  released codes fit our Hessians about 7\% \emph{worse} than plain
  round-to-nearest. The released codes are not an error-compensation of the
  public base; the weights themselves moved.
  \item \textbf{A public PTQ baseline is a calibration artifact}
  (\S\ref{sec:ptq}). We reproduce the published $2.106\times$ figure exactly on
  its fixed calibration passage, and show the same configurations collapse by
  $1{,}178\times$ to $\sim$$35{,}000\times$ on held-out windows with real
  calibration data.
\end{enumerate}

We also settle a mechanistic hypothesis: the large-$q$ mirror-descent /
regularized mirror-descent recipe~\citep{rmd,rmd2,proximalmf} is not the
mechanism behind the released weights (\S\ref{sec:recipe}). Direct tests inside
the ternary loop make it worse at every checkpoint, and the shape argument
explains why: a mirror map forms a single magnitude shell, while ternary needs
mass at both $0$ and $s$. Separately, we retract our own earlier 1.7B retention
figure after finding a holdout leak in our evaluation harness; the corrected,
clean value is reported with its range.

\paragraph{Scope and honesty.}
This is a forensic study of one artifact, not a benchmark suite and not a
replication of the recipe. Section~\ref{sec:limitations} lists the evidence
boundary and every threat to validity we know of, including the fact that some
numbers come from a 1.7B canary rather than the 27B, one corpus, and single
runs. Every falsified route is recorded in the failure register in
Appendix~\ref{app:failures} so that later work can re-assess it. No model
weights are redistributed; all analysis uses publicly released files and the
public base model.
